\chapter{Implement it}
\label{ch:implement}

``Implement it in NumPy'' is the interview question that separates people who have used deep
learning from people who understand it. This chapter assembles a tiny framework, every layer
with an explicit forward and backward pass, and proves each piece correct the strict way: the
NumPy model and a PyTorch model with the same initial weights are trained side by side on the
same batches, and their loss trajectories are asserted equal. Several pieces already appeared
where they were derived; the table points to them.

\begin{center}
	\footnotesize
	\begin{tabular}{p{0.42\linewidth}p{0.2\linewidth}p{0.28\linewidth}}
		\toprule
		component & where & checked against \\
		\midrule
		MLP backprop, gradient checking & \cref{ch:backprop}; here & autograd; finite differences \\
		reverse-mode autograd & \cref{ch:backprop} & autograd \\
		softmax regression & here & finite differences \\
		SGD, momentum, RMSProp, Adam, AdamW & \cref{ch:optim}; Adam again here & \pt{torch.optim} \\
		BatchNorm forward/backward, running statistics & \cref{ch:reg}; layer here & \pt{nn.BatchNorm1d} \\
		convolution forward/backward (im2col) & here & \pt{F.conv2d} + autograd \\
		LSTM cell; LSTM layer over a sequence & \cref{ch:rnn}; here & \pt{nn.LSTMCell}, \pt{nn.LSTM} \\
		scaled dot-product and multi-head attention & \cref{ch:transformer} & \pt{nn.MultiheadAttention} \\
		full encoder--decoder Transformer & \cref{ch:transformer} & PyTorch twin, same logits \\
		BPE tokeniser & \cref{ch:llm} & round trip \\
		character-level language model & here & PyTorch twin, same trajectory \\
		\bottomrule
	\end{tabular}
\end{center}

\section{A framework in sixty lines}
\prototype{\pt{Linear.backward} receives $\partial\Loss/\partial\text{output}$, stores
$\partial\Loss/\partial\W=g\T\x$ and returns $g\W$; everything else is bookkeeping.}

Each layer caches what its backward pass needs, exposes \pt{params} and, after \pt{backward}, \pt{grads}
with the same keys; the optimiser walks over both.

\codefile[\script{c12_implement}]{gen/dl/snip/c12_implement-framework.py}

\begin{example}[Softmax regression, gradient-checked]
	Softmax regression is \pt{Linear} followed by \pt{softmax\_ce}. On a 3-class, 10-feature dataset the analytic
	gradient of all 30 weights agrees with centred finite differences (step $10^{-6}$) to $10^{-7}$.
\end{example}

\section{BatchNorm as a layer}
\prototype{The same equations as \cref{ch:reg}, plus the two running averages and a \pt{training} flag.}

\codefile[\script{c12_implement}]{gen/dl/snip/c12_implement-batchnorm.py}

\begin{example}[A full training run, NumPy against PyTorch]
	A $10\to32\to3$ MLP with BatchNorm and ReLU, PyTorch's initial weights copied into the NumPy layers, both trained with
	Adam ($\eta=0.01$) on the same 50 random minibatches of 32. The two loss trajectories agree to $10^{-9}$ at every step
	(from $\gv{c12_implement/mlp_loss0}$ to $\gv{c12_implement/mlp_loss49}$); so do the final weights, BatchNorm's running
	variance, and the evaluation-mode outputs. Training accuracy ends at $\gv{c12_implement/mlp_acc}\%$. An identical
	trajectory over 50 Adam steps is a much stronger test than one gradient check: any error in a backward pass, in
	BatchNorm's unbiased running variance or in Adam's bias correction would make the curves drift apart.
\end{example}

\section{Convolution via im2col}
\prototype{Unfold every $k\times k\times C$ patch into a column; the convolution becomes one matrix product, and its
backward pass is two more.}

\vocab{im2col} turns a convolution into a matrix multiplication: the input $\shp{B,C,H,W}$ becomes columns
$\shp{B,Ck^2,H_oW_o}$, the weights a matrix $\shp{C_\text{out},Ck^2}$, and the output their product. Backward:
$\partial\Loss/\partial\W$ is the product of the output gradient with the columns, and $\partial\Loss/\partial\x$
multiplies by $\W\T$ and then scatter-adds the columns back to the image (\vocab{col2im}, the adjoint of im2col; overlapping
patches accumulate, exactly as in the transposed convolution of \cref{ch:cnn}).

\codefile[\script{c12_implement}]{gen/dl/snip/c12_implement-im2col.py}

For a batch of two $3\times7\times7$ inputs, four $3\times3$ filters, stride 2 and padding 1, the output has shape
$\gv{c12_implement/conv_out}$, and the output, $\partial\Loss/\partial\x$, $\partial\Loss/\partial\W$ and
$\partial\Loss/\partial\bb$ all equal PyTorch's. Production libraries use im2col-like lowering, FFT or Winograd
transforms, or direct kernels; the memory cost of im2col (each input pixel copied up to $k^2$ times) is the price of
reusing a fast matrix multiply.

\section{An LSTM layer}
\prototype{Loop the cell of \cref{ch:rnn} over time, carrying $(\h,\bc)$.}

\codefile[\script{c12_implement}]{gen/dl/snip/c12_implement-lstm_layer.py}

With weights taken from an \pt{nn.LSTM(4, 6)}, a 9-step sequence with batch 3 gives the same outputs and final cell state
as PyTorch. The backward pass is BPTT: run the cell's backward from $t=T$ down to $1$, adding the incoming
$\partial\Loss/\partial\h_t$ from the output at each step to the gradient arriving from step $t+1$, and accumulating the
weight gradients over time (gradient-check it as in \cref{ch:backprop}).

\section{A character-level language model}
\prototype{Three characters in, a distribution over the next character out: embedding, one $\tanh$ hidden layer, softmax.}

Two more layers complete the framework: an \pt{Embedding} whose backward pass scatter-adds gradients into the rows that
were used (\pt{np.add.at}, because a character can appear twice in one batch), and \pt{Tanh}.

\codefile[\script{c12_implement}]{gen/dl/snip/c12_implement-charlm.py}

\begin{example}[Training it, twice]
	A corpus of a few repeated sentences with a vocabulary of $\gv{c12_implement/clm_V}$ characters; context of 3,
	embeddings of size 8, 64 hidden units: $\gv{c12_implement/clm_params}$ parameters. Trained with the NumPy Adam above and,
	from the same initial weights and batches, with \pt{torch.optim.Adam}: the two 300-step loss trajectories agree to
	$10^{-8}$. The loss falls from $\gv{c12_implement/clm_loss0}$ (close to $\ln V=\gv{c12_implement/clm_lnV}$, as it
	should at initialisation) to $\gv{c12_implement/clm_loss_end}$, a per-character perplexity of
	$\gv{c12_implement/clm_ppl}$ on this tiny, repetitive corpus. Sampling from the NumPy model at temperature $0.5$, seeded
	with ``\pt{the}'':
	\outputfile{gen/dl/out/c12_implement-sample.txt}
	Locally fluent, globally aimless, and clearly memorised: a three-character window cannot track meaning. Longer context
	needs recurrence or attention (\cref{ch:rnn,ch:transformer}).
\end{example}

\begin{moral}
	Every layer is a forward function, a cache and a backward function that maps the output gradient to the input gradient
	and the parameter gradients. Get the shapes right and check against a reference, step by step.
\end{moral}

\section\problemhead

\begin{problem}\twochili
	Write the forward and backward pass of a \pt{Sigmoid} layer in the style of this chapter.
	\begin{hint}Cache the output.\end{hint}
	\begin{sol}
		\begin{pycode}
class Sigmoid:
    params = {}
    def forward(self, x):
        self.y = 1 / (1 + np.exp(-x))
        return self.y
    def backward(self, g):
        return g * self.y * (1 - self.y)
		\end{pycode}
		Caching $y$ rather than $x$ saves recomputing the exponential.
	\end{sol}
\end{problem}

\begin{problem}\twochili
	Why does \pt{Embedding.backward} use \pt{np.add.at} instead of \pt{gE[idx] += g}?
	\begin{hint}What if an index repeats?\end{hint}
	\begin{sol}
		With fancy indexing, \pt{gE[idx] += g} is a buffered assignment: when the same row appears several times in \pt{idx},
		only one contribution survives. \pt{np.add.at} performs an unbuffered accumulation, so every occurrence adds its gradient.
	\end{sol}
\end{problem}

\begin{problem}\twochili
	For a convolution with $C=64$ input channels, $k=3$, on a $56\times56$ map with stride 1 and padding 1, how many numbers does
	the im2col matrix hold per image, compared with the input?
	\begin{hint}$Ck^2\times H_oW_o$.\end{hint}
	\begin{sol}$64\cdot9\cdot56\cdot56=\gv{c12_implement/p_cols}$ versus $64\cdot56\cdot56=\gv{c12_implement/p_in}$: nine times
	larger, the cost of turning convolution into a matrix product.\end{sol}
\end{problem}
